\begin{center}
\textsc{Gauss} an \textsc{Bessel}.\quad Göttingen, 15.~November 1822.
\end{center}

\begingroup
\makeatletter
\renewcommand{\@makefnmark}{\hbox{\normalfont\textsuperscript{*})}}
\long\def\@makefntext#1{\parindent1.5em\noindent\textsuperscript{*})\ #1}
\makeatother
\par\noindent\hspace*{\parindent}.\ .\ .\ .\ .\ .\quad Die ausserordentlichen Schwierigkeiten, ein Dreiecksnetz in der Lüneburger Heide zu führen, kannte ich schon aus \textsc{Epaillys} Bericht, der es geradezu für unmöglich erklärt und seine Dreiecke, um den südlichen Theil von
\par\clearpage
\noindent352\hfill{\small\textbf{BRIEFWECHSEL MIT BESSEL.}}\par\medskip
\noindent Hannover mit Hamburg zu verbinden, über Bremen, die Weser herunter und so die Elbe wieder herauf geführt hat.
Und doch hatte er grosse Vortheile vor mir voraus; er beobachtete durchaus oben in seinen Signalthürmen, wo er sich viel leichter weitere Aussicht verschaffen konnte als ich, der überall zu ebener Erde gemessen hat; er brauchte beim Aushauen der Waldungen viel weniger auf Schonung Rücksicht zu nehmen; überall freier Transport, freie Arbeiter, etc., während ich alles ohne Ausnahme mit baarem Gelde bezahlen muss und oft die Kosten so sehr sehr weit über meine Vermuthung hinausgehen sehe; er hatte eine ganze Brigade von Ingenieurs zu Gehülfen, etc.
Diese Erwägungen liessen mich meine erste Recognoscirungsreise nicht ohne Ängstlichkeit vornehmen (gegen Ende Aprils).
Anfangs ging es selbst besser als ich erwartet hatte; ich fand, dass Garssen und Falkenberg sich unmittelbar mit Deister und Lichtenberg verbinden liessen; aber bei den Untersuchungen, wie von jenen beiden Punkten die Dreiecke weiter nördlich bis Lüneburg geführt werden könnten, fand ich das Terrain so widerspenstig, dass ich mehrere Male die Möglichkeit eines glücklichen Erfolgs bezweifelte und befürchtete, das ganze Unternehmen aufgeben zu müssen.
Das Land überall flach, keine dominirende Punkte, überall Holz, theils in grossen Waldungen, wie der Hassel, der Lüsing, das Becklinger Holz, etc., theils in unzählbaren kleinern Kämpen, die sich schachbrettartig vor einander schieben.
Die Versuche meines Gehülfen, auf der Westseite etwas brauchbares aufzufinden, waren ganz ohne Erfolg; mir selbst gelang es endlich nach den beschwerlichsten Versuchen, gleichsam im Herzen der Heide zwei Dreiecke 9.~10.~12 [Falkenberg-Hauselberg-Wulfsode] und 10.~12.~13 [Hauselberg-Wulfsode-Wilsede] zu etabliren; allein ich überzeugte mich zugleich, dass ich bei \textls[120]{dieser} Gattung von Arbeiten bald unterliegen würde, und setzte daher die weitere Aufsuchung einer Möglichkeit, diese Dreiecke mit den südlichern zu verknüpfen und weiter nördlich fortzuführen, auf die spätere Zeit hinaus, wo ich stärkeres Gehülfenpersonal (ich hatte nur einen Officier zur ersten Reise mitgenommen) und alle meine Instrumente bei mir haben würde.
Ich fing daher Mitte Junius die wirklichen Messungen in Lichtenberg an, und meine Hoffnung ist später auch so ziemlich erfüllt.
Da es ganz unmöglich befunden wurde, Garssen vermittelst Durchhaus mit Hauselberg zu verknüpfen, so wurde nach langem Suchen noch ein nördlicherer Punkt 14 [Scharnhorst]
\par\clearpage
\noindent{\small\textbf{ZUR HANNOVERSCHEN TRIANGULATION.}}\hfill353\par\medskip
\noindent gefunden, der mehr Hoffnung darbot.
Allein später zeigte sich dies absolut unmöglich, da der Wald zwischen 10 [Hauselberg] und 14 [Scharnhorst] auf zu hohem oder vielmehr nicht genug niedrigem Terrain lag, und ich musste mich glücklich schätzen, noch einen andern Punkt, Breithorn, zu finden, womit es endlich gelang.
Inzwischen war auch die Linie 9.~13 [Falkenberg-Wilsede] vermittelst eines sehr bedeutenden Durchhaus geöffnet und so hätte nun 10 [Hauselberg], wenn man 11 [Breithorn] früher gekannt hätte, ganz wegbleiben können.
Allein ich behielt jenen Punkt mit bei, und halte es für einen überaus schätzbaren Vortheil, dass in meinem Systeme drei Vierecke vorkommen, in denen alle sechs Richtungen wirklich hin und zurück gemessen sind.
Nachdem alle Winkel vorher in jedem Dreieck zur gehörigen Summe gehörig ausgeglichen waren, bedurfte es nur noch sehr kleiner Modificationen, meistens unter \(0^{\prime\prime}{,}1\), um diese drei Vierecke in volle Harmonie zu bringen.
Es wäre zu wünschen, dass man bei jeder Messung \textls[120]{solche} Prüfungen hätte.
Es gibt Messungen, wobei die Summen der drei Winkel überall zum Bewundern stimmen, und wo eine solche Prüfung zeigt, dass manche Winkel um \(2^{\prime\prime}\) bis \(3^{\prime\prime}\) gewiss unrichtig sind.
In der That ist die Prüfung vermittelst der Summe der Winkel à la portée von jedermann; die durch Diagonalen ist es weniger, so leicht sie auch für einen Mathematiker ist, und man kann sich der Vermuthung nicht erwehren, dass die erstere Prüfung zuweilen dazu gedient haben mag, wenn auch nicht die Beobachtungen zu verfälschen, doch etwas zu \textls[120]{wählen} (man bemerkt eine Tendenz dazu selbst bei \textsc{Delambre}).
Nördlich von Wulfsode ist zum Glück noch der Timpenberg gefunden, der unmittelbar mit Hamburg verbunden werden kann; der versuchte Durchhau von Timpenberg nach Lüneburg ist missglückt, weil das Land dazwischen zu wenig deprimirt war.
Die Richtungen von 10 nach 11 und von 11 nach 14 haben grosse Durchhaue erfordert.
Immer machte ich es mir zum Gesetz, mit der Rechnung allen Messungen, wie ich sie erhalten hatte, gleichen Schritt zu halten (bis auf die allerletzte Zeile), und nur dadurch ist es möglich geworden, alle Durchhaue mit der äussersten Präcision so durchzuführen, dass auch nicht Ein Stamm ohne Noth gefällt ist, oder die Unmöglichkeit der Durchhaue so früh wie möglich bestimmt zu erkennen.
So wusste ich z.~B.\ die Depression, unter der 14 [Scharnhorst] in 10 [Hauselberg] oder 20 [Lüneburg in 19 [Timpenberg] erscheinen musste, genau voraus; für den ersten
\par\smallskip\noindent{\footnotesize IX.}\hfill{\footnotesize45}\par
\clearpage
\noindent354\hfill{\small\textbf{BRIEFWECHSEL MIT BESSEL.}}\par\medskip
\noindent Fall war schon das Terrain vor dem Holz zu hoch, im zweiten, nachdem ein schmaler Spalt 2000 Schritt weit fortgeführt war.
Ich habe von Wilsede aus noch einen Punkt [Nindorf] 3000 Meter nördlich von Timpenberg festgelegt, der unmittelbar mit 20 [Lüneburg], 21 [Lauenburg], 23 [Hamburg] communicirt und wahrscheinlich mit 19 [Timpenberg] verknüpft werden kann, aber die ganze Strecke dahin muss durchgehauen werden.
Auf Timpenberg habe ich die Winkel zwischen 12 [Wulfsode], 13 [Wilsede], 23 [Hamburg] erst vorläufig gemessen, alle südlichern Punkte sind absolvirt sowie Wilsede.
Bei Scharnhorst bin ich zuletzt gewesen; dieser Punkt liesse sich vermittelst leichter Durchhaue mit Lichtenberg und Deister unmittelbar verbinden, wodurch Garssen entbehrlich würde; allein es wäre unmöglich gewesen, die Brauchbarkeit jenes Punkts zu Anfang auszumitteln.
Dieser Punkt wie mehrere andere sind ganz unscheinbare Plätze, von denen man gar nicht vermuthen sollte, dass sie so vielen Werth haben.
Bei vielen Richtungen, z.~B.\ von Lichtenberg nach Garssen und von Lichtenberg nach Falkenberg, geht die Linie so knapp über die Zwischenhindernisse, dass [der Zielpunkt] bei gewöhnlicher Refraction nur \textls[120]{wenige Secunden} sich darüber erhebt und nur zuweilen \(30^{\prime\prime}\) bis \(40^{\prime\prime}\) erreichte; von Falkenberg nach Wulfsode und vice versa kam das Licht bei schwächerer Refraction \textls[120]{gar nicht} herüber und hob sich immer erst in den spätern Nachmittagsstunden herauf.
Vom Lichtenberg selbst sah ich auf dem Falkenberg selten etwas und vice versa, das Heliotroplicht schwebte fast immer im freien Himmel (Distanz 85542 Meter); der grossen Distanz ungeachtet stimmen die Messungen ebenso gut oder vielmehr fast besser als auf ganz kleinen Distanzen von 10000 Meter Entfernung, wo das Heliotroplicht noch immer fast zu stark ist, wenngleich der Spiegel bis auf eine Öffnung von wenigen Quadratmillimetern bedeckt war.
Um an Zeit zu gewinnen, habe ich öfters auch selbst in Distanzen von 3 bis 4 bis 5 Meilen auf meine steinernen Postamente, \(3\frac{1}{2}\) Fuss hoch, selbst pointirt.
Dadurch sind die Messungen hin und wieder etwas weniger genau geworden, als wenn ich bloss Heliotroplicht gebraucht hätte, allein dann wäre ich in diesem Jahre lange nicht so weit gekommen.
Der grösste Fehler der Summe der drei Winkel war in diesem Jahre \(1^{\prime\prime}{,}76\), auch \textls[120]{überhaupt}\footnote{In 19 Dreiecken.} der grösste nächst dem
\par\clearpage
\noindent{\small\textbf{ZUR HANNOVERSCHEN TRIANGULATION.}}\hfill355\par\medskip
\noindent im Dreieck 4.~5.~6 [Hohehagen-Hils-Brocken], wo er \(3^{\prime\prime}{,}7\) beträgt und hauptsächlich dem schwierigen Pointiren auf den Brockenhaus-Thurm zuzuschreiben ist, welcher bei Sonnenschein nie gut geschnitten werden kann wegen der Phase.
Gern mässe ich die Winkel dieses Dreiecks, wenn es die Zeit erlaubte, noch einmal nach.
In diesem Jahre hatte ich oft drei Heliotrope zugleich in Activität, deren einen mein Sohn besorgte; es ist in der That ein prachtvoller Anblick.
Die Nachricht, welche ein gewisser \textsc{Schubach} über die Heliotrope im [Astronomischen] Jahrbuch 1825 gegeben hat, beruht auf einem Irrthum und hat gar nichts mit meinen Heliotropen gemein; diese sind künstliche Instrumente, deren Einrichtung mir erst \textls[120]{sehr viele} Mühe gekostet hat, die nun aber auch, wie ich glaube, nichts zu wünschen übrig lassen.
Ich habe von beiden Einrichtungen, die unter sich ganz verschieden sind, Exemplare für den General \textsc{Müffling} hier anfertigen lassen, auch \textsc{Gerling} hat zwei erhalten; vielleicht kann ich bald in \textsc{Schumachers} Astronomischen Nachrichten Abbildungen davon geben.
Zum Telegraphiren und um bei grossen Entfernungen den gegenüberstehenden Heliotropen erst die Richtung zu zeigen, brauchte ich oft einen grossen Spiegel von einem Fuss Quadrat, welcher wieder auf andere Art gelenkt wurde.
Der Anblick davon ist nach der Beschreibung meiner Gehülfen höchst prachtvoll gewesen; auf vier Meilen weit hat es dem blossen Auge zuweilen \textls[120]{wehe} gethan, lange hinzusehen. --- Doch jetzt genug hievon.

\par\noindent\hspace*{\parindent}Ihre Art, geodätische Beobachtungen zu behandeln, habe ich mit Vergnügen in \textsc{Schumachers} Astronomischen Nachrichten gesehen.
Sie wissen, dass dieser Gegenstand mich schon vor vielen Jahren beschäftigt hat.
Da Ihren Arbeiten nicht leicht etwas beigefügt werden kann, so würde, hätten unsere Wege sich begegnet, jene Bekanntmachung meine eigene Arbeit überflüssig gemacht haben.
Allein die Art, wie ich diesen Gegenstand behandelt habe, ist von der Ihrigen durch und durch verschieden, und so werde ich also in Zukunft bei Bekanntmachung meiner Messungen auch meine theoretischen Arbeiten ausführlich entwickeln.
Ich hoffe darin manches unerwartete geben zu können.
Aber diese Untersuchungen hängen mit einem reichen, fast unerschöpflich reichen Felde zusammen, und ich fühle oft mit inniger Wehmuth, bei dieser wie bei so vielen andern Gelegenheiten, wie meine äussern Verhältnisse mich an weitaussehenden theoretischen Arbeiten hindern.
Wenn solche ganz gedeihen sollen, muss man sich ihnen ganz hingeben können und nicht
\par\smallskip\noindent\hfill{\footnotesize45\textsuperscript{*}}\par
\clearpage
\noindent356\hfill{\small\textbf{BRIEFWECHSEL MIT BESSEL.}}\par\medskip
\noindent durch so heterogene Arbeiten wie Collegia lesen, alles kleinliche Detail beim Observiren und Rechnen der Beobachtungen, etc.\ etc.\ stündlich gehindert werden.\ .\ .\ .\ .\ .\ .

\par\noindent\hspace*{\parindent}Die Fatigen im heissen Sommer sind oft äusserst angreifend für mich gewesen, zuweilen so, dass ich glaubte, ich würde ihnen erliegen.
Auch das ist eine grosse Beschwerde bei den Arbeiten in der öden Lüneburger Heide, dass man öfters nur ein schlechtes Unterkommen und doch selbst ein solches nur meilenweit vom Arbeitspunkte haben kann.
Bei kühlem Wetter, welches meiner Constitution besser zusagt, befand ich mich im allgemeinen immer leidlich wohl, und jetzt kann ich über mein Befinden nicht klagen.\ .\ .\ .\ .\ .\ .

\par\noindent\hspace*{\parindent}Sie können aus obigem Bericht selbst sehen, was zur Vollendung meiner Triangulirung noch fehlt, aber einen eigentlichen Plan für nächsten Sommer kann ich jetzt noch nicht machen, es wird dabei auch vieles auf \textsc{Schumachers} Cooperation ankommen.
Es wäre \textls[120]{möglich}, dass ich, wenn ich mich bloss auf das Indispensable beschränke, schon im April und Mai die Dreiecke beendige und dann hieher zurückkomme; vielleicht im Julius hier und im August an einem nördlichern Punkte (Celle, Harburg, Hamburg?) mit dem Zenithsector messe.
Es könnte aber auch sein, wenn ich die Kosten nicht zu scheuen brauche, dass ich noch den ganzen Sommer auf die Triangulirung wende, um alles zur möglich grössten Vollkommenheit zu bringen.

\par\noindent\hspace*{\parindent}Vor einigen Tagen habe ich aus München ein Universalinstrument erhalten, das ich vor zwei Jahren vorzüglich behuf der Azimuthe bestellt hatte.
Inzwischen scheint nach meiner vorläufigen Reduction das von meinem Passageninstrument entlehnte und bis Hamburg übertragene Azimuth von dem, was mir \textsc{Schumacher} mitgetheilt hat (er hat auch vorläufig den Winkel zwischen Lüneburg und Wilsede, von wo ich ihm Heliotroplicht zusenden liess, gemessen), noch nicht \(1^{\prime\prime}{,}5\) zu differiren und also eine Azimuthmessung von meiner Seite an andern Punkten wohl ziemlich überflüssig zu sein.
Ob ich mich veranlasst sehen werde, mit dem Universalinstrument auch Polhöhen an mehrern Punkten zu messen, wird sich zeigen, wenn ich das Instrument erst eine Zeit lang gebraucht habe.
Es ist gut gearbeitet, scheint mir aber für den Gebrauch nicht recht bequem zu sein, besonders für ein kurzsichtiges Auge, welches zum Ablesen bei den Stellkreisen nicht gut zu kann.
Auf alle Fälle muss es sehr ermüdend sein, viel damit zu beobachten.
Meine Hori-
\par\clearpage
\noindent{\small\textbf{ZUR HANNOVERSCHEN TRIANGULATION.}}\hfill357\par\medskip
\noindent zontalwinkel habe ich alle mit einem 12-zölligen Theodolithen gemessen, der ein sehr vortreffliches Instrument ist.
Mit dem Zenithsector habe ich im vorigen Winter nur wenige Messungen gleichsam zur Probe gemacht; optische Kraft und Bequemlichkeit des Gebrauchs stehen dem \textsc{Reichenbach}schen Meridiankreis sehr nach.

\par\noindent\hspace*{\parindent}Meine Gradmessung als solche kann eigentlich für sich allein kein sehr wichtiges Resultat liefern; ich weiss aber nicht, wann die Vollendung der \textsc{Schumacher}schen zu hoffen ist.
So viel ich weiss, hat er seine Dreiecke erst etwa ein Drittel der ganzen Länge geführt, auch die Zenithsector-Messungen in Skagen sind bloss von seinem Gehülfen \textsc{Caroc} gemacht.
Die genaue Länge seiner Basis (circa 6000 Meter) kenne ich noch nicht; erst in diesem Herbst wollte er sie mit Hamburg verbinden, in seinem letzten Briefe erwähnt er aber gar nichts davon.
Ich könnte nun zwar selbst eine Basis messen und die Linie von Breithorn bis Scharnhorst scheint in ihrer ganzen Länge (11220 Meter) keine unübersteigliche Hindernisse darzubieten.
Allein auch abgesehen von den grossen Kosten gestehe ich, mich vor einer so höchst langweiligen Arbeit zu scheuen.
Meinen trigonometrischen Messungen habe ich immer eine interessante Seite abgewinnen können, da ihre tägliche Reduction immer einige Unterhaltung gab\footnote{Meine Beobachtungsmanier war, immer zu messen, was sich eben gut messen liess, ohne Rücksicht, ob es ein unmittelbarer Dreieckswinkel war. An manchen Stationen nahm ich einige Hülfspunkte, um auch dann nicht müssig zu sein, wenn nur Ein Heliotrop leuchtete. Die Franzosen haben, wie ich sehe, etwas ähnliches in Spanien gethan, aber die Messungen ganz unrichtig ausgeglichen.}.
Ich schnitt überdies auch alle sichtbaren Objecte bei Gelegenheit, um mich für die Landesgeographie nützlich zu beweisen, und ich muss sagen, dass ich dieses Geschäft mit seinen täglichen Ausgleichungen so lieb gewann, dass das Bemerken, Ausmitteln und Berechnen eines neuen Kirchthurms wohl ebenso viel Vergnügen machte, wie das Beobachten eines neuen Gestirns.
(Vor Gott ist's am Ende auch wohl einerlei, ob wir die Lage eines Kirchthurms auf einen Fuss oder die eines Sterns auf eine Secunde bestimmt haben.)
Allein bei einer Basismessung, sobald sie einmal im Gange ist, sehe ich für den Verstand auch gar nichts, was ihn reizen oder unterhalten könnte, und man muss sich bei einer vielleicht zwei Monate dauernden angestrengten täglichen Arbeit lediglich mit dem Gedanken auf-
\par\clearpage
\noindent358\hfill{\small\textbf{BRIEFWECHSEL MIT BESSEL.}}\par\medskip
\noindent recht halten, dass es eben doch einmal geschehen \textls[120]{muss}, um zuletzt eine Zahl zu haben.

\par\noindent\hspace*{\parindent}Die grosse Genauigkeit im Messen horizontaler Winkel durch Heliotroplicht auf die ungeheuersten Distanzen, und die Genauigkeit, womit man durch Universalinstrumente absolute Azimuthe messen kann, lassen mich glauben, dass man gegenwärtig eine Längengradmessung in schicklichem Terrain mit viel Vortheil ausführen könnte, indem man den Längenunterschied nicht auf Zeitbestimmung, sondern auf die Convergenz der Meridiane gründete.
Von den Pyrenäen sieht man den Mont Ventoux, von diesem die Alpen, von diesen bis Steiermark, von da bis Ungarn, etc.
Freilich wächst hiebei die Genauigkeit nicht wie beim Meridianbogen, wie die ganze Länge des Bogens (den eigentlich geodätischen Theil kann man dabei als fehlerfrei betrachten); also wenn man sich Stücke von ungefähr gleicher Grösse denkt, wächst beim Breitengrade die Genauigkeit wie die Anzahl der Stücke, beim Längengrade ebenfalls, wenn der Längenunterschied der Endpunkte astronomisch, z.~B.\ durch Sternbedeckungen, beobachtet wird; hingegen bei der Convergenz der Meridiane nur wie die Quadratwurzel der Zahl der Stücke (ungefähr); allein ich glaube doch, dass bei \textls[120]{sehr grossen} Stücken dies Verfahren mehr Genauigkeit gibt, als man z.~B.\ erhalten kann, wenn man den Längenunterschied durch Pulversignale vermittelst \textls[120]{Zeitbestimmung} sucht.
Am vortheilhaftesten wäre dies Verfahren in nördlichen Gegenden, wenn es dort sehr grosse Fernsichten gibt.
Sollten diese nicht in Norwegen und Schweden oder in einigen Gegenden von Russland zu finden sein?

\par\noindent\hspace*{\parindent}.\ .\ .\ .\ .\ .\quad Ich bin jetzt noch beim Ausgleichen, was ich leider einmal ganz umsonst gemacht habe, da durch ein Versehen von einer Station ein fehlerhaftes Tableau aufgenommen war.
Bei meiner Behandlung reagirt gewissermaassen jeder z.~B.\ in Wilsede gemessene Winkel auf alle übrigen bis Göttingen hin.
Ich werde Ihnen künftig einige Resultate anzeigen.
Nach vorläufiger Rechnung, Göttingen zu \(51^{0}\,31^{\prime}\,48^{\prime\prime}{,}7\) angenommen, fällt Hamburg in \(53^{0}\,33^{\prime}\,1^{\prime\prime}{,}76\), und \(0^{0}\,2^{\prime}\,2^{\prime\prime}{,}97\) östlich von Göttingen, das Absolute vorläufig auf \textsc{Zachs} Basis gestützt.
\par\noindent\hspace*{\parindent}.\ .\ .\ .\ .\ .
\endgroup
